\section{Conclusion}

Algorithmic redistricting produces substantially more competitive congressional
districts than enacted gerrymanders: 40\% vs.\ 17\% competitive on average
across five study states.
The efficiency gap narrows from $+5.7\%$ (Republican-favouring) to $-2.9\%$
(slight Democratic advantage reflecting geographic sorting).

The 23 pp increase in competitive districts is the largest single benefit
of algorithmic redistricting for democratic accountability.
More competitive districts hold representatives more accountable through the
general election, reduce the safe-seat dynamics that drive polarisation
(O.3), and may increase voter turnout (O.2).

The policy implication is direct: states that adopt algorithmic redistricting
will, as a byproduct of the process, gain substantially more competitive
congressional elections.
This competitive increase does not require any explicit partisan fairness
criterion in the algorithm --- it arises purely from the impossibility
defence (the algorithm cannot access partisan data) and the compactness
optimisation that the algorithm performs.
